% !TEX root = ../Main.tex
\chapter{平行关系}

平行分三层：线线平行、线面平行、面面平行。三层之间靠判定定理（往下走，把"线线"逐级推到"面面"）和性质定理（往上走，再逐级退回来）来回转换。下面把每一层常用的来源列清楚。

\section{直线之间（线线平行）}

\state{线线平行的来源}{
    \begin{enumerate}
        \item[(i)] 平面几何证法：平行四边形、平行线分线段成比例 $\Rightarrow$ 等角定理；
        \item[(ii)] 平行的传递性；
        \item[(iii)] $\bc l_1\parallel\alpha\\ \alpha\cap\beta=l_2\\ l_1\subset\beta\ec\ \Rightarrow\ l_1\parallel l_2$；
        \item[(iv)] $\bc \alpha\parallel\beta\\ \gamma\cap\alpha=l_1\\ \gamma\cap\beta=l_2\ec\ \Rightarrow\ l_1\parallel l_2$；
        \item[(v)] $\bc l_1\perp\alpha\\ l_2\perp\alpha\ec\ \Rightarrow\ l_1\parallel l_2$；
        \item[(vi)] $\bc l_1\perp\alpha\\ l_2\perp\beta\\ \alpha\parallel\beta\ec\ \Rightarrow\ l_1\parallel l_2$。
    \end{enumerate}
}

\section{线面平行}

\state{线面平行的来源}{
    \begin{enumerate}
        \item[(i)] $\bc l_1\parallel l_2\\ l_1\not\subset\alpha\\ l_2\subset\alpha\ec\ \Rightarrow\ l_1\parallel\alpha$（判定定理：面外的线平行于面内的线）；
        \item[(ii)] $\bc l_1\perp\beta\\ \alpha\perp\beta\\ l_1\not\subset\alpha\ec\ \Rightarrow\ l_1\parallel\alpha$；
        \item[(iii)] $\bc \alpha\parallel\beta\\ l\subset\alpha\ec\ \Rightarrow\ l\parallel\beta$。
    \end{enumerate}
}

\section{面面平行}

\state{面面平行的来源}{
    \begin{enumerate}
        \item[(i)] $\bc a\parallel\alpha,\ a\subset\beta\\ b\parallel\alpha,\ b\subset\beta\\ a\cap b=P\ec\ \Rightarrow\ \beta\parallel\alpha$（判定定理：一个面内两条相交直线都平行于另一个面）；
        \item[(ii)] $\bc l\perp\alpha\\ l\perp\beta\ec\ \Rightarrow\ \alpha\parallel\beta$。
    \end{enumerate}
}

三层之间的判定与性质，可以画成一张通路图：判定定理（右侧）从上往下把"线线"推到"面面"，性质定理（左侧）从下往上退回来。

\begin{center}
\begin{tikzpicture}[>=Stealth, node distance=1.5cm,
    box/.style={draw, rounded corners=1pt, minimum width=2.4cm, minimum height=0.8cm, font=\small}]
    \node[box] (ll) {线线平行};
    \node[box, below=of ll] (lm) {线面平行};
    \node[box, below=of lm] (mm) {面面平行};
    \draw[->] ($(lm.north)+(0.6,0)$) -- ($(ll.south)+(0.6,0)$) node[midway, right] {\scriptsize 性质};
    \draw[->] ($(ll.south)+(-0.6,0)$) -- ($(lm.north)+(-0.6,0)$) node[midway, left] {\scriptsize 判定};
    \draw[->] ($(mm.north)+(-0.6,0)$) -- ($(lm.south)+(-0.6,0)$) node[midway, left] {\scriptsize 判定};
    % 面面平行的性质可直接退到线线平行（外侧大箭头）
    \draw[->] (mm.east) -- ++(1.5,0) |- (ll.east) node[pos=0.25, right] {\scriptsize 性质};
\end{tikzpicture}
\end{center}
